% =========================================================
% RESUMEN
% El encabezado (\chapter*) está en main.tex.
% =========================================================

La inseguridad en México expone a la ciudadanía a riesgos latentes durante sus traslados cotidianos. Sin embargo, los peatones no siempre cuentan con herramientas para evaluar el nivel de vulnerabilidad de las rutas elegidas. Esta investigación propone el desarrollo de un prototipo de aplicación móvil enfocado en la generación de rutas peatonales seguras, fundamentado en algoritmos de búsqueda de caminos óptimos (\textit{Pathfinding} y \textit{Shortest Path Problem}), Teoría de Grafos, técnicas de cartografía participativa (\textit{crowdmapping}) y Sistemas de Información Geográfica (SIG). El modelo computacional pondera variables asociadas a la incidencia delictiva, morfología urbana, Criminología Ambiental y percepción ciudadana, con el fin de calcular trayectos con el menor nivel de exposición al riesgo.

\keywords{Aplicación móvil, Teoría de Grafos, Sistemas de Información Geográfica, Criminología Ambiental}
